% =============================================================================
% Clinical Decision Support — Consolidated LaTeX Templates
% =============================================================================
% This file contains all CDS templates in one place. Each template is wrapped
% in a section that can be extracted or compiled independently by copying it
% into its own .tex file with the shared preamble below.
%
% Sections:
%   1. Color Schemes
%   2. Biomarker Report Template
%   3. Clinical Pathway Template
%   4. Cohort Analysis Template
%   5. Treatment Recommendation Template
%
% Shared Preamble (use with any individual section):
%   \documentclass[10pt,letterpaper]{article}
%   \usepackage[margin=0.5in]{geometry}
%   \usepackage[utf8]{inputenc}
%   \usepackage[T1]{fontenc}
%   \usepackage{helvet} \renewcommand{\familydefault}{\sfdefault}
%   \usepackage{xcolor,tcolorbox,array,tabularx,booktabs,enumitem}
%   \usepackage{titlesec,fancyhdr,multicol,graphicx,float}
%   \usepackage{tikz} \usetikzlibrary{shapes,arrows,positioning,fit,calc}
%   % =============================================================================
% Clinical Decision Support — Consolidated LaTeX Templates
% =============================================================================
% This file contains all CDS templates in one place. Each template is wrapped
% in a section that can be extracted or compiled independently by copying it
% into its own .tex file with the shared preamble below.
%
% Sections:
%   1. Color Schemes
%   2. Biomarker Report Template
%   3. Clinical Pathway Template
%   4. Cohort Analysis Template
%   5. Treatment Recommendation Template
%
% Shared Preamble (use with any individual section):
%   \documentclass[10pt,letterpaper]{article}
%   \usepackage[margin=0.5in]{geometry}
%   \usepackage[utf8]{inputenc}
%   \usepackage[T1]{fontenc}
%   \usepackage{helvet} \renewcommand{\familydefault}{\sfdefault}
%   \usepackage{xcolor,tcolorbox,array,tabularx,booktabs,enumitem}
%   \usepackage{titlesec,fancyhdr,multicol,graphicx,float}
%   \usepackage{tikz} \usetikzlibrary{shapes,arrows,positioning,fit,calc}
%   % =============================================================================
% Clinical Decision Support — Consolidated LaTeX Templates
% =============================================================================
% This file contains all CDS templates in one place. Each template is wrapped
% in a section that can be extracted or compiled independently by copying it
% into its own .tex file with the shared preamble below.
%
% Sections:
%   1. Color Schemes
%   2. Biomarker Report Template
%   3. Clinical Pathway Template
%   4. Cohort Analysis Template
%   5. Treatment Recommendation Template
%
% Shared Preamble (use with any individual section):
%   \documentclass[10pt,letterpaper]{article}
%   \usepackage[margin=0.5in]{geometry}
%   \usepackage[utf8]{inputenc}
%   \usepackage[T1]{fontenc}
%   \usepackage{helvet} \renewcommand{\familydefault}{\sfdefault}
%   \usepackage{xcolor,tcolorbox,array,tabularx,booktabs,enumitem}
%   \usepackage{titlesec,fancyhdr,multicol,graphicx,float}
%   \usepackage{tikz} \usetikzlibrary{shapes,arrows,positioning,fit,calc}
%   % =============================================================================
% Clinical Decision Support — Consolidated LaTeX Templates
% =============================================================================
% This file contains all CDS templates in one place. Each template is wrapped
% in a section that can be extracted or compiled independently by copying it
% into its own .tex file with the shared preamble below.
%
% Sections:
%   1. Color Schemes
%   2. Biomarker Report Template
%   3. Clinical Pathway Template
%   4. Cohort Analysis Template
%   5. Treatment Recommendation Template
%
% Shared Preamble (use with any individual section):
%   \documentclass[10pt,letterpaper]{article}
%   \usepackage[margin=0.5in]{geometry}
%   \usepackage[utf8]{inputenc}
%   \usepackage[T1]{fontenc}
%   \usepackage{helvet} \renewcommand{\familydefault}{\sfdefault}
%   \usepackage{xcolor,tcolorbox,array,tabularx,booktabs,enumitem}
%   \usepackage{titlesec,fancyhdr,multicol,graphicx,float}
%   \usepackage{tikz} \usetikzlibrary{shapes,arrows,positioning,fit,calc}
%   \input{templates.tex}  % then \input only the section you need
% =============================================================================


% ╔═══════════════════════════════════════════════════════════════════════════╗
% ║  SECTION 1 — COLOR SCHEMES                                              ║
% ╚═══════════════════════════════════════════════════════════════════════════╝

% --- Primary Theme ---
\definecolor{headerblue}{RGB}{0,102,204}
\definecolor{highlightgray}{RGB}{240,240,240}

% --- Recommendation Strength ---
\definecolor{stronggreen}{RGB}{0,153,76}
\definecolor{strongdark}{RGB}{0,120,60}
\definecolor{conditionalyellow}{RGB}{255,193,7}
\definecolor{conditionalamber}{RGB}{255,160,0}
\definecolor{researchblue}{RGB}{33,150,243}
\definecolor{researchdark}{RGB}{25,118,210}
\definecolor{warningred}{RGB}{204,0,0}
\definecolor{dangerred}{RGB}{220,20,60}

% --- Urgency Levels ---
\definecolor{urgentred}{RGB}{220,20,60}
\definecolor{semiurgent}{RGB}{255,152,0}
\definecolor{routineblue}{RGB}{100,181,246}
\definecolor{actiongreen}{RGB}{0,153,76}

% --- Biomarker Categories ---
\definecolor{mutationred}{RGB}{244,67,54}
\definecolor{amplificationblue}{RGB}{33,150,243}
\definecolor{deletionpurple}{RGB}{156,39,176}
\definecolor{fusionpurple}{RGB}{156,39,176}
\definecolor{expressionorange}{RGB}{255,152,0}
\definecolor{tier1green}{RGB}{0,153,76}
\definecolor{tier2orange}{RGB}{255,152,0}
\definecolor{tier3gray}{RGB}{158,158,158}
\definecolor{biomarkerblue}{RGB}{51,102,204}

% --- Statistical Significance ---
\definecolor{significant}{RGB}{0,153,76}
\definecolor{trending}{RGB}{255,193,7}
\definecolor{nonsignificant}{RGB}{158,158,158}

% --- RECIST Response ---
\definecolor{completeresponse}{RGB}{0,153,76}
\definecolor{partialresponse}{RGB}{76,175,80}
\definecolor{stabledisease}{RGB}{255,193,7}
\definecolor{progressivedisease}{RGB}{244,67,54}

% --- Survival Outcomes ---
\definecolor{survivedgreen}{RGB}{0,153,76}
\definecolor{eventred}{RGB}{244,67,54}
\definecolor{censoredgray}{RGB}{158,158,158}

% --- CTCAE Adverse Event Severity ---
\definecolor{grade1}{RGB}{255,235,59}
\definecolor{grade2}{RGB}{255,193,7}
\definecolor{grade3}{RGB}{255,152,0}
\definecolor{grade4}{RGB}{244,67,54}
\definecolor{grade5}{RGB}{198,40,40}

% --- Colorblind-Safe (Okabe-Ito) ---
\definecolor{okabe1}{RGB}{230,159,0}
\definecolor{okabe2}{RGB}{86,180,233}
\definecolor{okabe3}{RGB}{0,158,115}
\definecolor{okabe4}{RGB}{240,228,66}
\definecolor{okabe5}{RGB}{0,114,178}
\definecolor{okabe6}{RGB}{213,94,0}
\definecolor{okabe7}{RGB}{204,121,167}

% Accessibility: always use text labels alongside color; test in grayscale;
% use Okabe-Ito palette for figures. Minimum contrast ratio 4.5:1.


% ╔═══════════════════════════════════════════════════════════════════════════╗
% ║  SECTION 2 — BIOMARKER REPORT TEMPLATE                                  ║
% ║  Genomic profiling report for molecular subtype classification.          ║
% ║  Standalone: \documentclass[10pt,letterpaper]{article} + shared preamble ║
% ╚═══════════════════════════════════════════════════════════════════════════╝
%
% Header/footer:
%   \fancyhead[L]{Genomic Profile Report: [PATIENT ID]}
%   \fancyfoot[C]{Confidential Laboratory Report - CLIA/CAP Certified}
%
% Structure:
%   Title → Patient/Specimen Info (tcolorbox) → Diagnosis → Results Summary
%   → Tier 1 (FDA-Approved) → Tier 2 (Investigational) → Tier 3 (VUS)
%   → Biomarkers Negative → TMB → MSI → Integrated Treatment Plan
%   → Clinical Trial Matching (table) → Methodology → Limitations
%   → Recommendations → References → Lab Director footer
%
% Key tcolorbox patterns:
%   Tier 1: colback=tier1green!5, colframe=tier1green
%     Mutation badge:  \colorbox{mutationred!60}{\textcolor{white}{\textbf{MUTATION}}}
%     Amplification:   \colorbox{amplificationblue!60}{...}
%   Tier 2: colback=tier2orange!5, colframe=tier2orange
%     Fusion badge:    \colorbox{fusionpurple!60}{...}
%   Tier 3: colback=tier3gray!10, colframe=tier3gray
%
% Clinical trial table columns: Trial | Intervention | Biomarker | Phase
% Methodology section: Assay info, variant classification (5-tier),
%   databases (OncoKB, CIViC, ClinVar, COSMIC)


% ╔═══════════════════════════════════════════════════════════════════════════╗
% ║  SECTION 3 — CLINICAL PATHWAY TEMPLATE                                  ║
% ║  TikZ decision-algorithm flowcharts for treatment sequencing.            ║
% ║  Standalone: \documentclass[10pt,letterpaper,landscape]{article}         ║
% ╚═══════════════════════════════════════════════════════════════════════════╝
%
% Header/footer:
%   \fancyhead[L]{Clinical Pathway: [CONDITION/DISEASE]}
%   \fancyfoot[C]{Evidence-Based Clinical Decision Pathway | For Professional Use Only}
%
% TikZ node styles:
\tikzstyle{cds@startstop} = [rectangle, rounded corners=8pt, minimum width=3cm,
  minimum height=1cm, text centered, draw=black, fill=headerblue!20,
  font=\small\bfseries]
\tikzstyle{cds@decision} = [diamond, minimum width=3cm, minimum height=1.2cm,
  text centered, draw=black, fill=conditionalyellow!40, font=\small,
  aspect=2, inner sep=0pt]
\tikzstyle{cds@process} = [rectangle, rounded corners=4pt, minimum width=3.5cm,
  minimum height=0.9cm, text centered, draw=black, fill=actiongreen!20,
  font=\small]
\tikzstyle{cds@urgent} = [rectangle, rounded corners=4pt, minimum width=3.5cm,
  minimum height=0.9cm, text centered, draw=urgentred, line width=1.5pt,
  fill=urgentred!15, font=\small\bfseries]
\tikzstyle{cds@routine} = [rectangle, rounded corners=4pt, minimum width=3.5cm,
  minimum height=0.9cm, text centered, draw=black, fill=routineblue!20,
  font=\small]
\tikzstyle{cds@info} = [rectangle, rounded corners=2pt, minimum width=2.5cm,
  minimum height=0.7cm, text centered, draw=researchblue, fill=researchblue!10,
  font=\footnotesize]
\tikzstyle{cds@arrow} = [thick,->,>=stealth]
\tikzstyle{cds@urgentarrow} = [ultra thick,->,>=stealth,color=urgentred]
%
% Flowchart structure:
%   Patient Presentation → Critical Criteria? --YES--> Immediate Action
%                                             --NO---> Continue Evaluation
%   Continue Evaluation → Risk Score >=X? --HIGH--> Admit ICU
%                                         --MOD---> Admit Observation
%                                         --LOW---> Outpatient
%   All branches → Definitive Management
%
% Legend box: shows Start/End, Decision Point, Action/Process symbols
%   and Urgency Color Coding (Urgent <1h, Semi-Urgent <24h, Routine >24h)
%
% Detailed Pathway Steps: 1) Initial Assessment, 2) Risk Stratification,
%   3) Treatment Initiation, 4) Monitoring and Reassessment
%
% Risk scoring table: Clinical Feature | Points → Risk Categories
% Evidence basis box: Key Supporting Evidence, Validation, Last Updated


% ╔═══════════════════════════════════════════════════════════════════════════╗
% ║  SECTION 4 — COHORT ANALYSIS TEMPLATE                                   ║
% ║  Biomarker-stratified patient cohort analysis with statistics.           ║
% ║  Standalone: \documentclass[10pt,letterpaper]{article} + shared preamble ║
% ╚═══════════════════════════════════════════════════════════════════════════╝
%
% Header/footer:
%   \fancyhead[L]{Clinical Decision Support: [COHORT NAME]}
%   \fancyfoot[C]{Confidential Medical Document - For Professional Use Only}
%
% Structure:
%   Title → Executive Summary (tcolorbox, highlightgray/headerblue)
%   → Cohort Characteristics: Demographics table (Group A vs B, p-value),
%     Biomarker Profile (biomarkerblue box)
%   → Treatment Exposures table
%   → Treatment Outcomes: Response Rates (RECIST v1.1), Survival (PFS, OS
%     with HR, CI, log-rank p), Kaplan-Meier figure placeholder
%   → Safety table (Any Grade + Grade 3-4, by group)
%   → Statistical Analysis: Methods, Multivariable Cox regression table
%   → Clinical Implications (highlightgreen box): Recommendations by group
%   → Subgroup Analyses: Interaction testing, exploratory subgroups
%   → Strengths and Limitations
%   → Conclusions, References, Footer


% ╔═══════════════════════════════════════════════════════════════════════════╗
% ║  SECTION 5 — TREATMENT RECOMMENDATION TEMPLATE                          ║
% ║  Evidence-based clinical practice guidelines with GRADE grading.         ║
% ║  Standalone: \documentclass[10pt,letterpaper]{article} + shared preamble ║
% ╚═══════════════════════════════════════════════════════════════════════════╝
%
% Header/footer:
%   \fancyhead[L]{Treatment Recommendations: [CONDITION]}
%   \fancyfoot[C]{Evidence-Based Clinical Guideline - For Professional Use Only}
%
% Recommendation Strength Legend:
%   STRONG (Grade 1): stronggreen!30 — benefits clearly outweigh risks
%   CONDITIONAL (Grade 2): conditionalyellow!30 — trade-offs, shared decision
%   RESEARCH (Grade R): researchblue!30 — insufficient evidence, trial preferred
%   Evidence Quality: A=High(RCTs), B=Moderate, C=Low(observational), D=Very low
%
% Structure:
%   Title → Strength Legend → Clinical Context (disease overview, population)
%   → Evidence Review: Key trials (design, population, endpoints, safety,
%     quality), Guideline concordance table (NCCN, ASCO, ESMO)
%   → Treatment Options:
%     First-Line: Option 1 (STRONG 1A, stronggreen box),
%                 Option 2 (CONDITIONAL 2B, conditionalyellow box),
%                 Option 3 (RESEARCH R, researchblue box)
%     Second-Line and Beyond
%   → Special Populations: Elderly (>=70), Renal impairment (eGFR table),
%     Hepatic impairment
%   → Clinical Decision Algorithm (TikZ flowchart)
%   → Monitoring Protocol: On-treatment table, Dose modification guidelines
%     (hematologic and non-hematologic toxicity)
%   → Recommendations by Clinical Scenario (GRADE boxes)
%   → Alternative Approaches: Not Recommended (warningred box)
%   → Supportive Care, Growth Factor Support
%   → Follow-Up and Surveillance table (Years 1-5+)
%   → Clinical Trial Opportunities
%   → Shared Decision-Making: Goals of care, decision aids (NNT)
%   → References, Committee footer
  % then \input only the section you need
% =============================================================================


% ╔═══════════════════════════════════════════════════════════════════════════╗
% ║  SECTION 1 — COLOR SCHEMES                                              ║
% ╚═══════════════════════════════════════════════════════════════════════════╝

% --- Primary Theme ---
\definecolor{headerblue}{RGB}{0,102,204}
\definecolor{highlightgray}{RGB}{240,240,240}

% --- Recommendation Strength ---
\definecolor{stronggreen}{RGB}{0,153,76}
\definecolor{strongdark}{RGB}{0,120,60}
\definecolor{conditionalyellow}{RGB}{255,193,7}
\definecolor{conditionalamber}{RGB}{255,160,0}
\definecolor{researchblue}{RGB}{33,150,243}
\definecolor{researchdark}{RGB}{25,118,210}
\definecolor{warningred}{RGB}{204,0,0}
\definecolor{dangerred}{RGB}{220,20,60}

% --- Urgency Levels ---
\definecolor{urgentred}{RGB}{220,20,60}
\definecolor{semiurgent}{RGB}{255,152,0}
\definecolor{routineblue}{RGB}{100,181,246}
\definecolor{actiongreen}{RGB}{0,153,76}

% --- Biomarker Categories ---
\definecolor{mutationred}{RGB}{244,67,54}
\definecolor{amplificationblue}{RGB}{33,150,243}
\definecolor{deletionpurple}{RGB}{156,39,176}
\definecolor{fusionpurple}{RGB}{156,39,176}
\definecolor{expressionorange}{RGB}{255,152,0}
\definecolor{tier1green}{RGB}{0,153,76}
\definecolor{tier2orange}{RGB}{255,152,0}
\definecolor{tier3gray}{RGB}{158,158,158}
\definecolor{biomarkerblue}{RGB}{51,102,204}

% --- Statistical Significance ---
\definecolor{significant}{RGB}{0,153,76}
\definecolor{trending}{RGB}{255,193,7}
\definecolor{nonsignificant}{RGB}{158,158,158}

% --- RECIST Response ---
\definecolor{completeresponse}{RGB}{0,153,76}
\definecolor{partialresponse}{RGB}{76,175,80}
\definecolor{stabledisease}{RGB}{255,193,7}
\definecolor{progressivedisease}{RGB}{244,67,54}

% --- Survival Outcomes ---
\definecolor{survivedgreen}{RGB}{0,153,76}
\definecolor{eventred}{RGB}{244,67,54}
\definecolor{censoredgray}{RGB}{158,158,158}

% --- CTCAE Adverse Event Severity ---
\definecolor{grade1}{RGB}{255,235,59}
\definecolor{grade2}{RGB}{255,193,7}
\definecolor{grade3}{RGB}{255,152,0}
\definecolor{grade4}{RGB}{244,67,54}
\definecolor{grade5}{RGB}{198,40,40}

% --- Colorblind-Safe (Okabe-Ito) ---
\definecolor{okabe1}{RGB}{230,159,0}
\definecolor{okabe2}{RGB}{86,180,233}
\definecolor{okabe3}{RGB}{0,158,115}
\definecolor{okabe4}{RGB}{240,228,66}
\definecolor{okabe5}{RGB}{0,114,178}
\definecolor{okabe6}{RGB}{213,94,0}
\definecolor{okabe7}{RGB}{204,121,167}

% Accessibility: always use text labels alongside color; test in grayscale;
% use Okabe-Ito palette for figures. Minimum contrast ratio 4.5:1.


% ╔═══════════════════════════════════════════════════════════════════════════╗
% ║  SECTION 2 — BIOMARKER REPORT TEMPLATE                                  ║
% ║  Genomic profiling report for molecular subtype classification.          ║
% ║  Standalone: \documentclass[10pt,letterpaper]{article} + shared preamble ║
% ╚═══════════════════════════════════════════════════════════════════════════╝
%
% Header/footer:
%   \fancyhead[L]{Genomic Profile Report: [PATIENT ID]}
%   \fancyfoot[C]{Confidential Laboratory Report - CLIA/CAP Certified}
%
% Structure:
%   Title → Patient/Specimen Info (tcolorbox) → Diagnosis → Results Summary
%   → Tier 1 (FDA-Approved) → Tier 2 (Investigational) → Tier 3 (VUS)
%   → Biomarkers Negative → TMB → MSI → Integrated Treatment Plan
%   → Clinical Trial Matching (table) → Methodology → Limitations
%   → Recommendations → References → Lab Director footer
%
% Key tcolorbox patterns:
%   Tier 1: colback=tier1green!5, colframe=tier1green
%     Mutation badge:  \colorbox{mutationred!60}{\textcolor{white}{\textbf{MUTATION}}}
%     Amplification:   \colorbox{amplificationblue!60}{...}
%   Tier 2: colback=tier2orange!5, colframe=tier2orange
%     Fusion badge:    \colorbox{fusionpurple!60}{...}
%   Tier 3: colback=tier3gray!10, colframe=tier3gray
%
% Clinical trial table columns: Trial | Intervention | Biomarker | Phase
% Methodology section: Assay info, variant classification (5-tier),
%   databases (OncoKB, CIViC, ClinVar, COSMIC)


% ╔═══════════════════════════════════════════════════════════════════════════╗
% ║  SECTION 3 — CLINICAL PATHWAY TEMPLATE                                  ║
% ║  TikZ decision-algorithm flowcharts for treatment sequencing.            ║
% ║  Standalone: \documentclass[10pt,letterpaper,landscape]{article}         ║
% ╚═══════════════════════════════════════════════════════════════════════════╝
%
% Header/footer:
%   \fancyhead[L]{Clinical Pathway: [CONDITION/DISEASE]}
%   \fancyfoot[C]{Evidence-Based Clinical Decision Pathway | For Professional Use Only}
%
% TikZ node styles:
\tikzstyle{cds@startstop} = [rectangle, rounded corners=8pt, minimum width=3cm,
  minimum height=1cm, text centered, draw=black, fill=headerblue!20,
  font=\small\bfseries]
\tikzstyle{cds@decision} = [diamond, minimum width=3cm, minimum height=1.2cm,
  text centered, draw=black, fill=conditionalyellow!40, font=\small,
  aspect=2, inner sep=0pt]
\tikzstyle{cds@process} = [rectangle, rounded corners=4pt, minimum width=3.5cm,
  minimum height=0.9cm, text centered, draw=black, fill=actiongreen!20,
  font=\small]
\tikzstyle{cds@urgent} = [rectangle, rounded corners=4pt, minimum width=3.5cm,
  minimum height=0.9cm, text centered, draw=urgentred, line width=1.5pt,
  fill=urgentred!15, font=\small\bfseries]
\tikzstyle{cds@routine} = [rectangle, rounded corners=4pt, minimum width=3.5cm,
  minimum height=0.9cm, text centered, draw=black, fill=routineblue!20,
  font=\small]
\tikzstyle{cds@info} = [rectangle, rounded corners=2pt, minimum width=2.5cm,
  minimum height=0.7cm, text centered, draw=researchblue, fill=researchblue!10,
  font=\footnotesize]
\tikzstyle{cds@arrow} = [thick,->,>=stealth]
\tikzstyle{cds@urgentarrow} = [ultra thick,->,>=stealth,color=urgentred]
%
% Flowchart structure:
%   Patient Presentation → Critical Criteria? --YES--> Immediate Action
%                                             --NO---> Continue Evaluation
%   Continue Evaluation → Risk Score >=X? --HIGH--> Admit ICU
%                                         --MOD---> Admit Observation
%                                         --LOW---> Outpatient
%   All branches → Definitive Management
%
% Legend box: shows Start/End, Decision Point, Action/Process symbols
%   and Urgency Color Coding (Urgent <1h, Semi-Urgent <24h, Routine >24h)
%
% Detailed Pathway Steps: 1) Initial Assessment, 2) Risk Stratification,
%   3) Treatment Initiation, 4) Monitoring and Reassessment
%
% Risk scoring table: Clinical Feature | Points → Risk Categories
% Evidence basis box: Key Supporting Evidence, Validation, Last Updated


% ╔═══════════════════════════════════════════════════════════════════════════╗
% ║  SECTION 4 — COHORT ANALYSIS TEMPLATE                                   ║
% ║  Biomarker-stratified patient cohort analysis with statistics.           ║
% ║  Standalone: \documentclass[10pt,letterpaper]{article} + shared preamble ║
% ╚═══════════════════════════════════════════════════════════════════════════╝
%
% Header/footer:
%   \fancyhead[L]{Clinical Decision Support: [COHORT NAME]}
%   \fancyfoot[C]{Confidential Medical Document - For Professional Use Only}
%
% Structure:
%   Title → Executive Summary (tcolorbox, highlightgray/headerblue)
%   → Cohort Characteristics: Demographics table (Group A vs B, p-value),
%     Biomarker Profile (biomarkerblue box)
%   → Treatment Exposures table
%   → Treatment Outcomes: Response Rates (RECIST v1.1), Survival (PFS, OS
%     with HR, CI, log-rank p), Kaplan-Meier figure placeholder
%   → Safety table (Any Grade + Grade 3-4, by group)
%   → Statistical Analysis: Methods, Multivariable Cox regression table
%   → Clinical Implications (highlightgreen box): Recommendations by group
%   → Subgroup Analyses: Interaction testing, exploratory subgroups
%   → Strengths and Limitations
%   → Conclusions, References, Footer


% ╔═══════════════════════════════════════════════════════════════════════════╗
% ║  SECTION 5 — TREATMENT RECOMMENDATION TEMPLATE                          ║
% ║  Evidence-based clinical practice guidelines with GRADE grading.         ║
% ║  Standalone: \documentclass[10pt,letterpaper]{article} + shared preamble ║
% ╚═══════════════════════════════════════════════════════════════════════════╝
%
% Header/footer:
%   \fancyhead[L]{Treatment Recommendations: [CONDITION]}
%   \fancyfoot[C]{Evidence-Based Clinical Guideline - For Professional Use Only}
%
% Recommendation Strength Legend:
%   STRONG (Grade 1): stronggreen!30 — benefits clearly outweigh risks
%   CONDITIONAL (Grade 2): conditionalyellow!30 — trade-offs, shared decision
%   RESEARCH (Grade R): researchblue!30 — insufficient evidence, trial preferred
%   Evidence Quality: A=High(RCTs), B=Moderate, C=Low(observational), D=Very low
%
% Structure:
%   Title → Strength Legend → Clinical Context (disease overview, population)
%   → Evidence Review: Key trials (design, population, endpoints, safety,
%     quality), Guideline concordance table (NCCN, ASCO, ESMO)
%   → Treatment Options:
%     First-Line: Option 1 (STRONG 1A, stronggreen box),
%                 Option 2 (CONDITIONAL 2B, conditionalyellow box),
%                 Option 3 (RESEARCH R, researchblue box)
%     Second-Line and Beyond
%   → Special Populations: Elderly (>=70), Renal impairment (eGFR table),
%     Hepatic impairment
%   → Clinical Decision Algorithm (TikZ flowchart)
%   → Monitoring Protocol: On-treatment table, Dose modification guidelines
%     (hematologic and non-hematologic toxicity)
%   → Recommendations by Clinical Scenario (GRADE boxes)
%   → Alternative Approaches: Not Recommended (warningred box)
%   → Supportive Care, Growth Factor Support
%   → Follow-Up and Surveillance table (Years 1-5+)
%   → Clinical Trial Opportunities
%   → Shared Decision-Making: Goals of care, decision aids (NNT)
%   → References, Committee footer
  % then \input only the section you need
% =============================================================================


% ╔═══════════════════════════════════════════════════════════════════════════╗
% ║  SECTION 1 — COLOR SCHEMES                                              ║
% ╚═══════════════════════════════════════════════════════════════════════════╝

% --- Primary Theme ---
\definecolor{headerblue}{RGB}{0,102,204}
\definecolor{highlightgray}{RGB}{240,240,240}

% --- Recommendation Strength ---
\definecolor{stronggreen}{RGB}{0,153,76}
\definecolor{strongdark}{RGB}{0,120,60}
\definecolor{conditionalyellow}{RGB}{255,193,7}
\definecolor{conditionalamber}{RGB}{255,160,0}
\definecolor{researchblue}{RGB}{33,150,243}
\definecolor{researchdark}{RGB}{25,118,210}
\definecolor{warningred}{RGB}{204,0,0}
\definecolor{dangerred}{RGB}{220,20,60}

% --- Urgency Levels ---
\definecolor{urgentred}{RGB}{220,20,60}
\definecolor{semiurgent}{RGB}{255,152,0}
\definecolor{routineblue}{RGB}{100,181,246}
\definecolor{actiongreen}{RGB}{0,153,76}

% --- Biomarker Categories ---
\definecolor{mutationred}{RGB}{244,67,54}
\definecolor{amplificationblue}{RGB}{33,150,243}
\definecolor{deletionpurple}{RGB}{156,39,176}
\definecolor{fusionpurple}{RGB}{156,39,176}
\definecolor{expressionorange}{RGB}{255,152,0}
\definecolor{tier1green}{RGB}{0,153,76}
\definecolor{tier2orange}{RGB}{255,152,0}
\definecolor{tier3gray}{RGB}{158,158,158}
\definecolor{biomarkerblue}{RGB}{51,102,204}

% --- Statistical Significance ---
\definecolor{significant}{RGB}{0,153,76}
\definecolor{trending}{RGB}{255,193,7}
\definecolor{nonsignificant}{RGB}{158,158,158}

% --- RECIST Response ---
\definecolor{completeresponse}{RGB}{0,153,76}
\definecolor{partialresponse}{RGB}{76,175,80}
\definecolor{stabledisease}{RGB}{255,193,7}
\definecolor{progressivedisease}{RGB}{244,67,54}

% --- Survival Outcomes ---
\definecolor{survivedgreen}{RGB}{0,153,76}
\definecolor{eventred}{RGB}{244,67,54}
\definecolor{censoredgray}{RGB}{158,158,158}

% --- CTCAE Adverse Event Severity ---
\definecolor{grade1}{RGB}{255,235,59}
\definecolor{grade2}{RGB}{255,193,7}
\definecolor{grade3}{RGB}{255,152,0}
\definecolor{grade4}{RGB}{244,67,54}
\definecolor{grade5}{RGB}{198,40,40}

% --- Colorblind-Safe (Okabe-Ito) ---
\definecolor{okabe1}{RGB}{230,159,0}
\definecolor{okabe2}{RGB}{86,180,233}
\definecolor{okabe3}{RGB}{0,158,115}
\definecolor{okabe4}{RGB}{240,228,66}
\definecolor{okabe5}{RGB}{0,114,178}
\definecolor{okabe6}{RGB}{213,94,0}
\definecolor{okabe7}{RGB}{204,121,167}

% Accessibility: always use text labels alongside color; test in grayscale;
% use Okabe-Ito palette for figures. Minimum contrast ratio 4.5:1.


% ╔═══════════════════════════════════════════════════════════════════════════╗
% ║  SECTION 2 — BIOMARKER REPORT TEMPLATE                                  ║
% ║  Genomic profiling report for molecular subtype classification.          ║
% ║  Standalone: \documentclass[10pt,letterpaper]{article} + shared preamble ║
% ╚═══════════════════════════════════════════════════════════════════════════╝
%
% Header/footer:
%   \fancyhead[L]{Genomic Profile Report: [PATIENT ID]}
%   \fancyfoot[C]{Confidential Laboratory Report - CLIA/CAP Certified}
%
% Structure:
%   Title → Patient/Specimen Info (tcolorbox) → Diagnosis → Results Summary
%   → Tier 1 (FDA-Approved) → Tier 2 (Investigational) → Tier 3 (VUS)
%   → Biomarkers Negative → TMB → MSI → Integrated Treatment Plan
%   → Clinical Trial Matching (table) → Methodology → Limitations
%   → Recommendations → References → Lab Director footer
%
% Key tcolorbox patterns:
%   Tier 1: colback=tier1green!5, colframe=tier1green
%     Mutation badge:  \colorbox{mutationred!60}{\textcolor{white}{\textbf{MUTATION}}}
%     Amplification:   \colorbox{amplificationblue!60}{...}
%   Tier 2: colback=tier2orange!5, colframe=tier2orange
%     Fusion badge:    \colorbox{fusionpurple!60}{...}
%   Tier 3: colback=tier3gray!10, colframe=tier3gray
%
% Clinical trial table columns: Trial | Intervention | Biomarker | Phase
% Methodology section: Assay info, variant classification (5-tier),
%   databases (OncoKB, CIViC, ClinVar, COSMIC)


% ╔═══════════════════════════════════════════════════════════════════════════╗
% ║  SECTION 3 — CLINICAL PATHWAY TEMPLATE                                  ║
% ║  TikZ decision-algorithm flowcharts for treatment sequencing.            ║
% ║  Standalone: \documentclass[10pt,letterpaper,landscape]{article}         ║
% ╚═══════════════════════════════════════════════════════════════════════════╝
%
% Header/footer:
%   \fancyhead[L]{Clinical Pathway: [CONDITION/DISEASE]}
%   \fancyfoot[C]{Evidence-Based Clinical Decision Pathway | For Professional Use Only}
%
% TikZ node styles:
\tikzstyle{cds@startstop} = [rectangle, rounded corners=8pt, minimum width=3cm,
  minimum height=1cm, text centered, draw=black, fill=headerblue!20,
  font=\small\bfseries]
\tikzstyle{cds@decision} = [diamond, minimum width=3cm, minimum height=1.2cm,
  text centered, draw=black, fill=conditionalyellow!40, font=\small,
  aspect=2, inner sep=0pt]
\tikzstyle{cds@process} = [rectangle, rounded corners=4pt, minimum width=3.5cm,
  minimum height=0.9cm, text centered, draw=black, fill=actiongreen!20,
  font=\small]
\tikzstyle{cds@urgent} = [rectangle, rounded corners=4pt, minimum width=3.5cm,
  minimum height=0.9cm, text centered, draw=urgentred, line width=1.5pt,
  fill=urgentred!15, font=\small\bfseries]
\tikzstyle{cds@routine} = [rectangle, rounded corners=4pt, minimum width=3.5cm,
  minimum height=0.9cm, text centered, draw=black, fill=routineblue!20,
  font=\small]
\tikzstyle{cds@info} = [rectangle, rounded corners=2pt, minimum width=2.5cm,
  minimum height=0.7cm, text centered, draw=researchblue, fill=researchblue!10,
  font=\footnotesize]
\tikzstyle{cds@arrow} = [thick,->,>=stealth]
\tikzstyle{cds@urgentarrow} = [ultra thick,->,>=stealth,color=urgentred]
%
% Flowchart structure:
%   Patient Presentation → Critical Criteria? --YES--> Immediate Action
%                                             --NO---> Continue Evaluation
%   Continue Evaluation → Risk Score >=X? --HIGH--> Admit ICU
%                                         --MOD---> Admit Observation
%                                         --LOW---> Outpatient
%   All branches → Definitive Management
%
% Legend box: shows Start/End, Decision Point, Action/Process symbols
%   and Urgency Color Coding (Urgent <1h, Semi-Urgent <24h, Routine >24h)
%
% Detailed Pathway Steps: 1) Initial Assessment, 2) Risk Stratification,
%   3) Treatment Initiation, 4) Monitoring and Reassessment
%
% Risk scoring table: Clinical Feature | Points → Risk Categories
% Evidence basis box: Key Supporting Evidence, Validation, Last Updated


% ╔═══════════════════════════════════════════════════════════════════════════╗
% ║  SECTION 4 — COHORT ANALYSIS TEMPLATE                                   ║
% ║  Biomarker-stratified patient cohort analysis with statistics.           ║
% ║  Standalone: \documentclass[10pt,letterpaper]{article} + shared preamble ║
% ╚═══════════════════════════════════════════════════════════════════════════╝
%
% Header/footer:
%   \fancyhead[L]{Clinical Decision Support: [COHORT NAME]}
%   \fancyfoot[C]{Confidential Medical Document - For Professional Use Only}
%
% Structure:
%   Title → Executive Summary (tcolorbox, highlightgray/headerblue)
%   → Cohort Characteristics: Demographics table (Group A vs B, p-value),
%     Biomarker Profile (biomarkerblue box)
%   → Treatment Exposures table
%   → Treatment Outcomes: Response Rates (RECIST v1.1), Survival (PFS, OS
%     with HR, CI, log-rank p), Kaplan-Meier figure placeholder
%   → Safety table (Any Grade + Grade 3-4, by group)
%   → Statistical Analysis: Methods, Multivariable Cox regression table
%   → Clinical Implications (highlightgreen box): Recommendations by group
%   → Subgroup Analyses: Interaction testing, exploratory subgroups
%   → Strengths and Limitations
%   → Conclusions, References, Footer


% ╔═══════════════════════════════════════════════════════════════════════════╗
% ║  SECTION 5 — TREATMENT RECOMMENDATION TEMPLATE                          ║
% ║  Evidence-based clinical practice guidelines with GRADE grading.         ║
% ║  Standalone: \documentclass[10pt,letterpaper]{article} + shared preamble ║
% ╚═══════════════════════════════════════════════════════════════════════════╝
%
% Header/footer:
%   \fancyhead[L]{Treatment Recommendations: [CONDITION]}
%   \fancyfoot[C]{Evidence-Based Clinical Guideline - For Professional Use Only}
%
% Recommendation Strength Legend:
%   STRONG (Grade 1): stronggreen!30 — benefits clearly outweigh risks
%   CONDITIONAL (Grade 2): conditionalyellow!30 — trade-offs, shared decision
%   RESEARCH (Grade R): researchblue!30 — insufficient evidence, trial preferred
%   Evidence Quality: A=High(RCTs), B=Moderate, C=Low(observational), D=Very low
%
% Structure:
%   Title → Strength Legend → Clinical Context (disease overview, population)
%   → Evidence Review: Key trials (design, population, endpoints, safety,
%     quality), Guideline concordance table (NCCN, ASCO, ESMO)
%   → Treatment Options:
%     First-Line: Option 1 (STRONG 1A, stronggreen box),
%                 Option 2 (CONDITIONAL 2B, conditionalyellow box),
%                 Option 3 (RESEARCH R, researchblue box)
%     Second-Line and Beyond
%   → Special Populations: Elderly (>=70), Renal impairment (eGFR table),
%     Hepatic impairment
%   → Clinical Decision Algorithm (TikZ flowchart)
%   → Monitoring Protocol: On-treatment table, Dose modification guidelines
%     (hematologic and non-hematologic toxicity)
%   → Recommendations by Clinical Scenario (GRADE boxes)
%   → Alternative Approaches: Not Recommended (warningred box)
%   → Supportive Care, Growth Factor Support
%   → Follow-Up and Surveillance table (Years 1-5+)
%   → Clinical Trial Opportunities
%   → Shared Decision-Making: Goals of care, decision aids (NNT)
%   → References, Committee footer
  % then \input only the section you need
% =============================================================================


% ╔═══════════════════════════════════════════════════════════════════════════╗
% ║  SECTION 1 — COLOR SCHEMES                                              ║
% ╚═══════════════════════════════════════════════════════════════════════════╝

% --- Primary Theme ---
\definecolor{headerblue}{RGB}{0,102,204}
\definecolor{highlightgray}{RGB}{240,240,240}

% --- Recommendation Strength ---
\definecolor{stronggreen}{RGB}{0,153,76}
\definecolor{strongdark}{RGB}{0,120,60}
\definecolor{conditionalyellow}{RGB}{255,193,7}
\definecolor{conditionalamber}{RGB}{255,160,0}
\definecolor{researchblue}{RGB}{33,150,243}
\definecolor{researchdark}{RGB}{25,118,210}
\definecolor{warningred}{RGB}{204,0,0}
\definecolor{dangerred}{RGB}{220,20,60}

% --- Urgency Levels ---
\definecolor{urgentred}{RGB}{220,20,60}
\definecolor{semiurgent}{RGB}{255,152,0}
\definecolor{routineblue}{RGB}{100,181,246}
\definecolor{actiongreen}{RGB}{0,153,76}

% --- Biomarker Categories ---
\definecolor{mutationred}{RGB}{244,67,54}
\definecolor{amplificationblue}{RGB}{33,150,243}
\definecolor{deletionpurple}{RGB}{156,39,176}
\definecolor{fusionpurple}{RGB}{156,39,176}
\definecolor{expressionorange}{RGB}{255,152,0}
\definecolor{tier1green}{RGB}{0,153,76}
\definecolor{tier2orange}{RGB}{255,152,0}
\definecolor{tier3gray}{RGB}{158,158,158}
\definecolor{biomarkerblue}{RGB}{51,102,204}

% --- Statistical Significance ---
\definecolor{significant}{RGB}{0,153,76}
\definecolor{trending}{RGB}{255,193,7}
\definecolor{nonsignificant}{RGB}{158,158,158}

% --- RECIST Response ---
\definecolor{completeresponse}{RGB}{0,153,76}
\definecolor{partialresponse}{RGB}{76,175,80}
\definecolor{stabledisease}{RGB}{255,193,7}
\definecolor{progressivedisease}{RGB}{244,67,54}

% --- Survival Outcomes ---
\definecolor{survivedgreen}{RGB}{0,153,76}
\definecolor{eventred}{RGB}{244,67,54}
\definecolor{censoredgray}{RGB}{158,158,158}

% --- CTCAE Adverse Event Severity ---
\definecolor{grade1}{RGB}{255,235,59}
\definecolor{grade2}{RGB}{255,193,7}
\definecolor{grade3}{RGB}{255,152,0}
\definecolor{grade4}{RGB}{244,67,54}
\definecolor{grade5}{RGB}{198,40,40}

% --- Colorblind-Safe (Okabe-Ito) ---
\definecolor{okabe1}{RGB}{230,159,0}
\definecolor{okabe2}{RGB}{86,180,233}
\definecolor{okabe3}{RGB}{0,158,115}
\definecolor{okabe4}{RGB}{240,228,66}
\definecolor{okabe5}{RGB}{0,114,178}
\definecolor{okabe6}{RGB}{213,94,0}
\definecolor{okabe7}{RGB}{204,121,167}

% Accessibility: always use text labels alongside color; test in grayscale;
% use Okabe-Ito palette for figures. Minimum contrast ratio 4.5:1.


% ╔═══════════════════════════════════════════════════════════════════════════╗
% ║  SECTION 2 — BIOMARKER REPORT TEMPLATE                                  ║
% ║  Genomic profiling report for molecular subtype classification.          ║
% ║  Standalone: \documentclass[10pt,letterpaper]{article} + shared preamble ║
% ╚═══════════════════════════════════════════════════════════════════════════╝
%
% Header/footer:
%   \fancyhead[L]{Genomic Profile Report: [PATIENT ID]}
%   \fancyfoot[C]{Confidential Laboratory Report - CLIA/CAP Certified}
%
% Structure:
%   Title → Patient/Specimen Info (tcolorbox) → Diagnosis → Results Summary
%   → Tier 1 (FDA-Approved) → Tier 2 (Investigational) → Tier 3 (VUS)
%   → Biomarkers Negative → TMB → MSI → Integrated Treatment Plan
%   → Clinical Trial Matching (table) → Methodology → Limitations
%   → Recommendations → References → Lab Director footer
%
% Key tcolorbox patterns:
%   Tier 1: colback=tier1green!5, colframe=tier1green
%     Mutation badge:  \colorbox{mutationred!60}{\textcolor{white}{\textbf{MUTATION}}}
%     Amplification:   \colorbox{amplificationblue!60}{...}
%   Tier 2: colback=tier2orange!5, colframe=tier2orange
%     Fusion badge:    \colorbox{fusionpurple!60}{...}
%   Tier 3: colback=tier3gray!10, colframe=tier3gray
%
% Clinical trial table columns: Trial | Intervention | Biomarker | Phase
% Methodology section: Assay info, variant classification (5-tier),
%   databases (OncoKB, CIViC, ClinVar, COSMIC)


% ╔═══════════════════════════════════════════════════════════════════════════╗
% ║  SECTION 3 — CLINICAL PATHWAY TEMPLATE                                  ║
% ║  TikZ decision-algorithm flowcharts for treatment sequencing.            ║
% ║  Standalone: \documentclass[10pt,letterpaper,landscape]{article}         ║
% ╚═══════════════════════════════════════════════════════════════════════════╝
%
% Header/footer:
%   \fancyhead[L]{Clinical Pathway: [CONDITION/DISEASE]}
%   \fancyfoot[C]{Evidence-Based Clinical Decision Pathway | For Professional Use Only}
%
% TikZ node styles:
\tikzstyle{cds@startstop} = [rectangle, rounded corners=8pt, minimum width=3cm,
  minimum height=1cm, text centered, draw=black, fill=headerblue!20,
  font=\small\bfseries]
\tikzstyle{cds@decision} = [diamond, minimum width=3cm, minimum height=1.2cm,
  text centered, draw=black, fill=conditionalyellow!40, font=\small,
  aspect=2, inner sep=0pt]
\tikzstyle{cds@process} = [rectangle, rounded corners=4pt, minimum width=3.5cm,
  minimum height=0.9cm, text centered, draw=black, fill=actiongreen!20,
  font=\small]
\tikzstyle{cds@urgent} = [rectangle, rounded corners=4pt, minimum width=3.5cm,
  minimum height=0.9cm, text centered, draw=urgentred, line width=1.5pt,
  fill=urgentred!15, font=\small\bfseries]
\tikzstyle{cds@routine} = [rectangle, rounded corners=4pt, minimum width=3.5cm,
  minimum height=0.9cm, text centered, draw=black, fill=routineblue!20,
  font=\small]
\tikzstyle{cds@info} = [rectangle, rounded corners=2pt, minimum width=2.5cm,
  minimum height=0.7cm, text centered, draw=researchblue, fill=researchblue!10,
  font=\footnotesize]
\tikzstyle{cds@arrow} = [thick,->,>=stealth]
\tikzstyle{cds@urgentarrow} = [ultra thick,->,>=stealth,color=urgentred]
%
% Flowchart structure:
%   Patient Presentation → Critical Criteria? --YES--> Immediate Action
%                                             --NO---> Continue Evaluation
%   Continue Evaluation → Risk Score >=X? --HIGH--> Admit ICU
%                                         --MOD---> Admit Observation
%                                         --LOW---> Outpatient
%   All branches → Definitive Management
%
% Legend box: shows Start/End, Decision Point, Action/Process symbols
%   and Urgency Color Coding (Urgent <1h, Semi-Urgent <24h, Routine >24h)
%
% Detailed Pathway Steps: 1) Initial Assessment, 2) Risk Stratification,
%   3) Treatment Initiation, 4) Monitoring and Reassessment
%
% Risk scoring table: Clinical Feature | Points → Risk Categories
% Evidence basis box: Key Supporting Evidence, Validation, Last Updated


% ╔═══════════════════════════════════════════════════════════════════════════╗
% ║  SECTION 4 — COHORT ANALYSIS TEMPLATE                                   ║
% ║  Biomarker-stratified patient cohort analysis with statistics.           ║
% ║  Standalone: \documentclass[10pt,letterpaper]{article} + shared preamble ║
% ╚═══════════════════════════════════════════════════════════════════════════╝
%
% Header/footer:
%   \fancyhead[L]{Clinical Decision Support: [COHORT NAME]}
%   \fancyfoot[C]{Confidential Medical Document - For Professional Use Only}
%
% Structure:
%   Title → Executive Summary (tcolorbox, highlightgray/headerblue)
%   → Cohort Characteristics: Demographics table (Group A vs B, p-value),
%     Biomarker Profile (biomarkerblue box)
%   → Treatment Exposures table
%   → Treatment Outcomes: Response Rates (RECIST v1.1), Survival (PFS, OS
%     with HR, CI, log-rank p), Kaplan-Meier figure placeholder
%   → Safety table (Any Grade + Grade 3-4, by group)
%   → Statistical Analysis: Methods, Multivariable Cox regression table
%   → Clinical Implications (highlightgreen box): Recommendations by group
%   → Subgroup Analyses: Interaction testing, exploratory subgroups
%   → Strengths and Limitations
%   → Conclusions, References, Footer


% ╔═══════════════════════════════════════════════════════════════════════════╗
% ║  SECTION 5 — TREATMENT RECOMMENDATION TEMPLATE                          ║
% ║  Evidence-based clinical practice guidelines with GRADE grading.         ║
% ║  Standalone: \documentclass[10pt,letterpaper]{article} + shared preamble ║
% ╚═══════════════════════════════════════════════════════════════════════════╝
%
% Header/footer:
%   \fancyhead[L]{Treatment Recommendations: [CONDITION]}
%   \fancyfoot[C]{Evidence-Based Clinical Guideline - For Professional Use Only}
%
% Recommendation Strength Legend:
%   STRONG (Grade 1): stronggreen!30 — benefits clearly outweigh risks
%   CONDITIONAL (Grade 2): conditionalyellow!30 — trade-offs, shared decision
%   RESEARCH (Grade R): researchblue!30 — insufficient evidence, trial preferred
%   Evidence Quality: A=High(RCTs), B=Moderate, C=Low(observational), D=Very low
%
% Structure:
%   Title → Strength Legend → Clinical Context (disease overview, population)
%   → Evidence Review: Key trials (design, population, endpoints, safety,
%     quality), Guideline concordance table (NCCN, ASCO, ESMO)
%   → Treatment Options:
%     First-Line: Option 1 (STRONG 1A, stronggreen box),
%                 Option 2 (CONDITIONAL 2B, conditionalyellow box),
%                 Option 3 (RESEARCH R, researchblue box)
%     Second-Line and Beyond
%   → Special Populations: Elderly (>=70), Renal impairment (eGFR table),
%     Hepatic impairment
%   → Clinical Decision Algorithm (TikZ flowchart)
%   → Monitoring Protocol: On-treatment table, Dose modification guidelines
%     (hematologic and non-hematologic toxicity)
%   → Recommendations by Clinical Scenario (GRADE boxes)
%   → Alternative Approaches: Not Recommended (warningred box)
%   → Supportive Care, Growth Factor Support
%   → Follow-Up and Surveillance table (Years 1-5+)
%   → Clinical Trial Opportunities
%   → Shared Decision-Making: Goals of care, decision aids (NNT)
%   → References, Committee footer
